\documentclass[11pt]{article}
\usepackage[a4paper,margin=1in]{geometry}
\usepackage{fontspec}
\usepackage[boldfont=false]{xeCJK}
\setCJKmainfont{FandolSong-Regular.otf}
\usepackage{amsmath}
\usepackage{amssymb}
\usepackage{array}
\usepackage{booktabs}
\usepackage{graphicx}
\usepackage{adjustbox}
\usepackage{enumitem}
\setlist[itemize]{topsep=2pt,partopsep=0pt,parsep=0pt,itemsep=2pt,leftmargin=1.4em}
\setlist[enumerate]{topsep=2pt,partopsep=0pt,parsep=0pt,itemsep=2pt,leftmargin=1.6em}
\usepackage{parskip}
\usepackage{fvextra}
\fvset{breaklines=true,breakanywhere=true,fontsize=\small}
\setlength{\parindent}{0pt}

\begin{document}

\begin{center}
  {\LARGE\bfseries Polymerisation}\\[0.35em]
  {\large A Level Chemistry}
\end{center}
\vspace{0.6em}

\section*{Condensation polymerisation}

In \textbf{condensation polymerisation} 缩合聚合, monomers join into a long chain \textbf{and} a small molecule (such as water or $\text{HCl}$) is lost each time a new bond forms. Each monomer needs \textbf{two} reactive groups, so the chain can grow at both ends.

\subsection*{Polyesters}

A \textbf{polyester} 聚酯 has many ester links along its chain. You can make one from:

\begin{itemize}
\item a \textbf{diol} 二醇 (two $\text{–OH}$ groups) and a \textbf{dicarboxylic acid} 二羧酸 (two $\text{–COOH}$ groups), or a dioyl chloride.

\item a single hydroxycarboxylic acid, which has both an $\text{–OH}$ and a $\text{–COOH}$.

\end{itemize}

\subsection*{Polyamides}

A \textbf{polyamide} 聚酰胺 has many amide links. You can make one from:

\begin{itemize}
\item a \textbf{diamine} 二胺 (two $\text{–NH}_2$ groups) and a dicarboxylic acid or a dioyl chloride.

\item a single aminocarboxylic acid, or from \textbf{amino acids} 氨基酸 joining together (proteins are natural polyamides).

\end{itemize}

\subsection*{Repeat units and monomers}

The \textbf{repeat unit} 重复单元 of a condensation polymer contains parts of \textbf{both} monomers, minus the atoms lost as the small molecule. To find the \textbf{monomers} 单体 from a section of polymer, break the chain at each ester or amide link, then add back $\text{–OH}$ and $\text{–H}$ (or $\text{–Cl}$).

\section*{Predicting the type of polymerisation}

\begin{center}
\adjustbox{max width=\linewidth}{%
\begin{tabular}{ll}
\toprule
\textbf{Clue} & \textbf{Type} \\
\midrule
the monomer has a C=C double bond, and \textbf{nothing} else is lost & \textbf{addition polymerisation} 加成聚合 \\
each monomer has \textbf{two} functional groups, and a small molecule is lost; the chain has ester or amide links & condensation polymerisation \\
\bottomrule
\end{tabular}}
\end{center}

\section*{Degradable polymers}

\begin{itemize}
\item \textbf{poly(alkene)s} 聚烯烃 are chemically inert — they have only strong, non-polar C–C and C–H bonds, so they are hard to \textbf{biodegrade} 可生物降解 and last a long time.

\item some polymers are made so that light can break them down (they are photodegradable).

\item polyesters and polyamides \textbf{are} biodegradable, because their ester and amide links can be broken by acidic or alkaline \textbf{hydrolysis} 水解.

\end{itemize}


\end{document}
